\section{Broad-Family Support, Profiles, and Influence}\label{app:families}

\subsection{Family-by-quintile support}

Table~\ref{tab:app_family_support} reports, for every SOC two-digit family, the number of retained occupations, preperiod stock share, exposure range, quintiles present, direct-tail overlap, and nuisance-adjusted Q5 information share.  Only arts/media (SOC27), health practitioners (SOC29), health support (SOC31), and sales (SOC41) span Q1 and Q5 under the reconstructed comparison.  Other families connect the Q5--Q1 regression coefficient through intermediate quintiles and the imposed common Q2--Q5 slopes.

\begin{center}
\footnotesize
\begin{longtable}{@{}r p{0.31\textwidth} r r r l c@{}}
\caption{Broad-family exposure support}\label{tab:app_family_support}\\
\toprule
SOC2 & Family & Occs. & Stock (\%) & Beta range & Quintiles & Q1 and Q5 \\ 
\midrule
\endfirsthead
\multicolumn{7}{l}{\footnotesize\itshape Table \thetable{} continued}\\
\toprule
SOC2 & Family & Occs. & Stock (\%) & Beta range & Quintiles & Q1 and Q5 \\ 
\midrule
\endhead
\midrule\multicolumn{7}{r}{\footnotesize Continued on next page}\\
\endfoot
\bottomrule
\endlastfoot
11 & Management & 29 & 13.56 & 0.367--0.586 & 3|4|5 & No \\
13 & Business and Financial Operations & 28 & 6.22 & 0.405--0.636 & 3|4|5 & No \\
15 & Computer and Mathematical & 15 & 4.31 & 0.500--0.950 & 4|5 & No \\
17 & Architecture and Engineering & 19 & 2.28 & 0.406--0.727 & 3|4|5 & No \\
19 & Life, Physical, and Social Science & 21 & 1.25 & 0.344--0.723 & 3|4|5 & No \\
21 & Community and Social Service & 12 & 1.31 & 0.111--0.439 & 1|2|3 & No \\
23 & Legal & 4 & 1.34 & 0.425--0.571 & 3|4|5 & No \\
25 & Educational Instruction and Library & 7 & 4.53 & 0.229--0.567 & 2|3|5 & No \\
27 & Arts, Design, Entertainment, Sports, and Media & 25 & 1.80 & 0.000--0.958 & 1|2|3|4|5 & Yes \\
29 & Healthcare Practitioners and Technical & 37 & 6.38 & 0.067--0.618 & 1|2|3|4|5 & Yes \\
31 & Healthcare Support & 14 & 3.66 & 0.074--0.875 & 1|2|3|5 & Yes \\
33 & Protective Service & 17 & 2.28 & 0.036--0.415 & 1|2|3 & No \\
35 & Food Preparation and Serving Related & 10 & 3.27 & 0.000--0.511 & 1|2|4 & No \\
37 & Building and Grounds Cleaning and Maintenance & 6 & 2.51 & 0.000--0.355 & 1|2|3 & No \\
39 & Personal Care and Service & 17 & 2.51 & 0.015--0.400 & 1|2|3 & No \\
41 & Sales and Related & 17 & 10.06 & 0.133--0.806 & 1|2|3|4|5 & Yes \\
43 & Office and Administrative Support & 47 & 11.21 & 0.216--0.975 & 2|3|4|5 & No \\
45 & Farming, Fishing, and Forestry & 5 & 0.15 & 0.043--0.292 & 1|2 & No \\
47 & Construction and Extraction & 30 & 6.14 & 0.000--0.417 & 1|2|3 & No \\
49 & Installation, Maintenance, and Repair & 31 & 3.56 & 0.000--0.410 & 1|2|3 & No \\
51 & Production & 52 & 4.14 & 0.000--0.438 & 1|2|3 & No \\
53 & Transportation and Material Moving & 25 & 7.52 & 0.031--0.413 & 1|2|3 & No \\
\end{longtable}
\tabnote{Stock is each family's share of preperiod employment on the 468-occupation rebuilt support.  Beta ranges and quintile memberships use the corrected 2017--2022 preperiod treatment contract.  Only rows marked Yes enter the changed-population direct-tail exercise.}
\end{center}


Table~\ref{tab:app_family_information} makes the loss and concentration of identifying information explicit.  Family-by-month conditioning retains 29.7 percent of the pooled Q5 target's nuisance-adjusted information, while the direct-tail comparison has only 6.2 effective occupations and its five largest occupation contributions supply 77.8 percent of information.  Its normal-theory MDE$_{80}$ is 0.4576 log point.  The continuous within-family companion is more diffuse and has MDE$_{80}=0.0313$ per one within-family standard deviation, but it imposes a common slope across families.

\begin{table}[!htbp]
\centering
\caption{Family-support information and precision diagnostics}\label{tab:app_family_information}
\resizebox{\textwidth}{!}{%
\begin{tabular}{lrrrrrrr}
\toprule
Target and structure & Occ. & $I$ (millions) & $I/I_{\mathrm{raw}}$ & $I/I_{\mathrm{parent}}$ & Effective occ. & Top-five share & Normal MDE$_{80}$ \\
\midrule
Q5--Q1, pooled & 468 & 26.289 & 0.378 & 1.000 & 43.2 & 0.245 & 0.1266 \\
Q5--Q1, SOC2 $\times$ month & 468 & 7.803 & 0.112 & 0.297 & 53.3 & 0.218 & 0.1998 \\
Direct-tail Q5--Q1, SOC2 $\times$ month & 29 & 0.669 & 0.080 & 0.308 & 6.2 & 0.778 & 0.4576 \\
Continuous within-SOC2 slope, SOC2 $\times$ month & 468 & 261.278 & 0.583 & 0.965 & 48.0 & 0.249 & 0.0313 \\
\bottomrule
\end{tabular}}
\tabnote{$I$ is nuisance-adjusted Fisher information after projecting the target score on the model's nuisance score; $I_{\mathrm{raw}}$ is the raw weighted treatment sum of squares.  ``Parent'' is the corresponding pooled model on the same target and population, so the direct-tail and continuous rows each use their own pooled parent.  Because target scales and populations differ, absolute $I$ is not comparable across all rows.  Effective occupations are the inverse Herfindahl count of occupation-level information contributions; the top-five share is the sum of the five largest contributions.  Fitted information and its effective count are design-concentration diagnostics: neither replaces the nominal occupation-cluster count nor measures first-stage CPS sampling precision.  MDE$_{80}=2.80\widehat{SE}$ uses the occupation-cluster standard error and is a normal-theory precision diagnostic, not a rejection result.  The direct-tail MDE is 0.4576 log point; the continuous MDE is 0.0313 per one within-family standard deviation (0.1084 raw beta units).}
\end{table}


The direct-tail population retains only occupations in the four spanning families that belong to Q1 or Q5.  It contains 29 occupations and 5.03 percent of primary preperiod stock.  In the SOC2-by-calendar-month model its coefficient is 0.1494, with interval $[-0.1694,0.4682]$.  This benchmark is valuable because it exposes direct overlap; it is not a preferred estimate for the original population, and the interval includes large effects of either sign.

Table~\ref{tab:app_direct_tail_occupations} names every occupation supplying
that comparison.  The imbalance is economically visible: for example, the
health-support Q5 endpoint consists only of medical transcriptionists, while
its Q1 endpoint includes nursing assistants and home health aides.  Naming the
membership clarifies what the direct-tail estimate compares; it does not make
the sparse endpoints representative of their broad families.

\begin{center}
\scriptsize
\begin{longtable}{@{}r c r p{0.47\textwidth} r@{}}
\caption{Occupations supplying direct within-family tail support}\label{tab:app_direct_tail_occupations}\\
\toprule
SOC2 & Tail & \shortstack{Census\\occupation} & Occupation & \shortstack{Preperiod stock\\(million person-months)} \\
\midrule
\endfirsthead
\multicolumn{5}{l}{\footnotesize\itshape Table \thetable{} continued}\\
\toprule
SOC2 & Tail & \shortstack{Census\\occupation} & Occupation & \shortstack{Preperiod stock\\(million person-months)} \\
\midrule
\endhead
\midrule\multicolumn{5}{r}{\footnotesize Continued on next page}\\
\endfoot
\bottomrule
\endlastfoot
27 & Q1 & 2721 & Athletes and sports competitors & 1.83 \\
27 & Q1 & 2740 & Dancers and choreographers & 0.94 \\
27 & Q5 & 2850 & Writers and authors & 14.91 \\
27 & Q5 & 2825 & Public relations specialists & 8.61 \\
27 & Q5 & 2830 & Editors & 8.44 \\
27 & Q5 & 2861 & Interpreters and translators & 6.14 \\
27 & Q5 & 2810 & News analysts, reporters, and journalists & 5.41 \\
27 & Q5 & 2840 & Technical writers & 4.38 \\
27 & Q5 & 2751 & Music directors and composers & 2.74 \\
27 & Q5 & 2805 & Broadcast announcers and radio disc jockeys & 2.19 \\
27 & Q5 & 2862 & Court reporters and simultaneous captioners & 2.10 \\
29 & Q1 & 3423 & Surgical technologists & 6.77 \\
29 & Q1 & 3402 & Paramedics & 6.47 \\
29 & Q1 & 3422 & Psychiatric technicians & 3.59 \\
29 & Q5 & 3515 & Medical records specialists & 11.77 \\
29 & Q5 & 3030 & Dietitians and nutritionists & 7.71 \\
31 & Q1 & 3603 & Nursing assistants & 86.89 \\
31 & Q1 & 3601 & Home health aides & 31.36 \\
31 & Q1 & 3630 & Massage therapists & 10.70 \\
31 & Q1 & 3605 & Orderlies and psychiatric aides & 3.48 \\
31 & Q5 & 3646 & Medical transcriptionists & 2.62 \\
41 & Q1 & 4900 & Models, demonstrators, and product promoters & 2.18 \\
41 & Q5 & 4850 & Sales representatives, wholesale and manufacturing & 81.01 \\
41 & Q5 & 4810 & Insurance sales agents & 38.93 \\
41 & Q5 & 4840 & Sales representatives of services, except advertising, insurance, financial services, and travel & 36.99 \\
41 & Q5 & 4820 & Securities, commodities, and financial services sales agents & 15.73 \\
41 & Q5 & 4800 & Advertising sales agents & 11.43 \\
41 & Q5 & 4830 & Travel agents & 4.26 \\
41 & Q5 & 4940 & Telemarketers & 2.87 \\
\end{longtable}
\tabnote{This is the complete 29-occupation membership of the changed-population direct-tail exercise, ordered by family, tail, and stock.  Preperiod stock is the cumulative CPS-weighted employed-person stock across the 71 observed months from January 2017 through November 2022: for each occupation, the producer sums the monthly young (ages 22--25) and older (ages 26--65) employment stocks.  The displayed unit is one million employed-person months, not average-month employment.  The list makes the direct comparisons auditable; it is not a ranking by estimated effect or influence.}
\end{center}


\subsection{Supported edges and broader beta support}

The fully interacted within-family model yields 89 supported pairwise
quintile-edge functionals.  All 89 estimates and their standard errors are
retained in the machine-readable results.  They are distinct family-specific
comparisons rather than repeated national Q5--Q1 estimates.  Their joint
covariance has rank 50, so an 89-degree-of-freedom Wald test is invalid and is
blocked in the results rather than replaced by a selected list of pointwise
$p$-values.  The associated 50-dimensional family-specific coefficient system
is reported separately with rank-aware inference.

Table~\ref{tab:app_pairwise_aggregates} gives a transparent summary of the
supported edges.  Each row averages only family-specific contrasts for which
both endpoint quintiles are observed, using fixed preperiod endpoint-stock
weights.  Because the eligible family set changes across rows, these
aggregates do not impose a common national quintile profile and cannot be read
as ten estimates of one population parameter.

\begin{table}[!htbp]
\centering
\caption{Aggregates of supported within-family quintile contrasts}\label{tab:app_pairwise_aggregates}
\begin{tabular}{lrrrr}
\toprule
Contrast & Eligible families & Estimate & Standard error & Normal $p$-value \\
\midrule
Q2--Q1 & 14 & $-0.0095$ & 0.0368 & .796 \\
Q3--Q1 & 12 & $\phantom{-}0.0132$ & 0.0541 & .807 \\
Q4--Q1 & 4 & $-0.0618$ & 0.0577 & .284 \\
Q5--Q1 & 4 & $\phantom{-}0.1020$ & 0.0781 & .192 \\
Q3--Q2 & 14 & $\phantom{-}0.0702$ & 0.0348 & .044 \\
Q4--Q2 & 5 & $-0.0033$ & 0.0354 & .926 \\
Q5--Q2 & 6 & $\phantom{-}0.1177$ & 0.0498 & .018 \\
Q4--Q3 & 9 & $-0.0841$ & 0.0458 & .066 \\
Q5--Q3 & 11 & $-0.0446$ & 0.0965 & .644 \\
Q5--Q4 & 10 & $-0.1082$ & 0.0485 & .026 \\
\bottomrule
\end{tabular}
\tabnote{For each pair, only families containing both listed quintiles enter.  Family-specific edge coefficients are aggregated using fixed pair-endpoint preperiod employment-stock weights, and the covariance retains the corresponding linear functional of the full coefficient system.  The ten rows have different eligible-family populations and are not a monotone national dose-response sequence.  Normal $p$-values are descriptive and were not used to select rows; all 89 underlying supported family-specific edges and their simultaneous intervals remain machine readable.}
\end{table}


Table~\ref{tab:app_broader_support} asks whether sparse tails arise solely from
requiring Webb software exposure.  Relaxing that requirement expands the
population from 468 to 490 occupations, but the added occupations represent
only 1.88 percent of broader-support preperiod stock.  Direct tails expand from
29 to 62 occupations and from 5.03 to 12.61 percent of stock because SOC43
joins the four original tail-spanning families.  The overlap therefore improves
at the margin but remains confined to five of 23 families.  Coefficients change
little in this descriptive population expansion; this is not a bound and does
not make the two populations equivalent.

\begin{table}[!htbp]
\centering
\caption{Supported tails when Webb availability is relaxed}\label{tab:app_broader_support}
\resizebox{\textwidth}{!}{%
\begin{tabular}{lrrrrrr}
\toprule
Population and classification & Occupations & Families & Tail-spanning families & Tail occupations & Tail stock (\%) & Pooled coefficient \\
\midrule
Primary support, with Webb & 468 & 22 & 4 & 29 & 5.03 & $-0.1321$ \\
Primary support, Webb omitted & 468 & 22 & 4 & 29 & 5.03 & $-0.1338$ \\
Broader beta-valid support, primary cuts & 490 & 23 & 5 & 62 & 12.61 & $-0.1344$ \\
Broader beta-valid support, recomputed cuts & 490 & 23 & 5 & 62 & 12.61 & $-0.1348$ \\
\bottomrule
\end{tabular}}
\tabnote{The primary support requires both beta exposure and Webb software exposure; omitting Webb from the regression on that support does not change membership.  The broader support requires beta exposure but not Webb and adds 22 occupations, equal to 1.88 percent of broader-support preperiod stock.  With the primary cuts or recomputed cuts it adds office and administrative support (SOC43) as the fifth family containing both Q1 and Q5.  ``Tail occupations'' and ``tail stock'' count Q1 and Q5 occupations within tail-spanning families.  The coefficient change from 468 to 490 occupations changes the population and is descriptive, not a paired test.  On fixed support, the Webb-omitted minus Webb-included coefficient difference is $-0.0017$, with paired interval $[-0.0102,0.0069]$; nondetection is not equivalence.}
\end{table}


\subsection{Full exposure profile}

The pooled Q2--Q5 coefficients are jointly different from zero
($\chi^2_4=11.1383$, asymptotic $p=.0251$, multiplier $p=.0212$), but pairwise
movements do not establish a monotone dose response.  After SOC2-by-calendar-
month conditioning, the four coefficients are jointly indistinguishable from
zero ($\chi^2_4=2.0708$, asymptotic $p=.7227$, multiplier $p=.7422$).  The
four-coefficient paired movement likewise does not reject
($\chi^2_4=5.4169$, asymptotic $p=.2471$, multiplier $p=.2368$), despite the
single Q5 movement's pointwise interval excluding zero.  Table~\ref{tab:app_profile}
reports pointwise and simultaneous intervals and the common-draw movement for
every group.

\begin{table}[!htbp]
\centering
\caption{Exposure profile before and after broad-family monthly paths}\label{tab:app_profile}
\scriptsize
\resizebox{\textwidth}{!}{%
\begin{tabular}{lrrrrrrr}
\toprule
& \multicolumn{3}{c}{Pooled} & \multicolumn{3}{c}{SOC2 $\times$ month} & Paired movement \\ 
\cmidrule(lr){2-4}\cmidrule(lr){5-7}
Group & Coef. & Pointwise 95\% CI & Simultaneous 95\% CI & Coef. & Pointwise 95\% CI & Simultaneous 95\% CI & Coef. [95\% CI] \\ 
\midrule
Q2 & -0.0784 & [-0.1546,-0.0023] & [-0.1716,0.0147] & -0.0247 & [-0.1228,0.0734] & [-0.1431,0.0937] & 0.0537 [-0.0051,0.1126] \\
Q3 & -0.0500 & [-0.1224,0.0225] & [-0.1383,0.0383] & 0.0169 & [-0.0859,0.1197] & [-0.1043,0.1381] & 0.0669 [-0.0047,0.1386] \\
Q4 & -0.0896 & [-0.1578,-0.0214] & [-0.1737,-0.0055] & -0.0126 & [-0.1302,0.1051] & [-0.1510,0.1259] & 0.0770 [-0.0138,0.1679] \\
Q5 & -0.1321 & [-0.2203,-0.0439] & [-0.2393,-0.0249] & -0.0217 & [-0.1595,0.1162] & [-0.1836,0.1403] & 0.1104 [0.0085,0.2124] \\
\bottomrule
\end{tabular}}
\tabnote{Q1 is omitted.  Simultaneous intervals use each model's common four-target maximum-$|t|$ critical value.  Paired movements are SOC2-by-calendar-month minus pooled coefficients on common occupation draws.  All rows use the rebuilt corrected-preperiod treatment contract and 9,999 draws.}
\end{table}


The continuous companion demeans beta within family using preperiod employment weights and imposes a common linear slope.  With SOC2-by-calendar-month conditioning, the rebuilt coefficient is $-0.0025$ per 0.1084 raw beta units.  Its interval is $[-0.0240,0.0191]$ in that standardized within-family unit.  The unconditioned continuous coefficient is essentially zero as well.  This is a different functional form from the tail contrast.

\subsection{Exact numerator, denominator, and family accounting}

Table~\ref{tab:app_stock_accounting} separates the model-free aggregate
Q5--Q1 stock-ratio movement into its young numerator and older denominator.
The older tail difference reinforces the negative young movement but is
smaller in absolute magnitude.  A second exact accounting separates changes
in within-family young-to-older ratios from changes in older family weights.
Neither accounting is forced to equal the nonlinear regression target, and no
component is interpreted as a causal percentage explained.

\begin{table}[!htbp]
\centering
\caption{Exact descriptive accounting for the Q5--Q1 stock-ratio change}\label{tab:app_stock_accounting}
\begin{tabular}{lrr}
\toprule
\multicolumn{3}{l}{\textit{Panel A. Underlying equal-month employment stocks}} \\
Series & Preperiod & Postperiod \\
\midrule
Q1 young, ages 22--25 (millions) & 2.5902 & 2.9044 \\
Q1 older, ages 26--65 (millions) & 21.1411 & 21.1983 \\
Q5 young, ages 22--25 (millions) & 1.9562 & 1.9654 \\
Q5 older, ages 26--65 (millions) & 21.1303 & 22.2367 \\
Q1 young-to-older ratio & 0.1225 & 0.1370 \\
Q5 young-to-older ratio & 0.0926 & 0.0884 \\
\addlinespace
\multicolumn{3}{l}{\textit{Panel B. Exact Q5--Q1 changes and decomposition}} \\
Component & \multicolumn{2}{r}{Log points} \\
\midrule
Young-stock Q5--Q1 change & \multicolumn{2}{r}{$-0.1098$} \\
Older-stock Q5--Q1 change & \multicolumn{2}{r}{$\phantom{-}0.0483$} \\
Young-relative change (young minus older) & \multicolumn{2}{r}{$-0.1581$} \\
\addlinespace
Within-family young-to-older ratio component & \multicolumn{2}{r}{$-0.1293$} \\
Changing older-family-weight component & \multicolumn{2}{r}{$-0.0288$} \\
Observed boundary-mass component & \multicolumn{2}{r}{$\phantom{-}0.0000$} \\
\bottomrule
\end{tabular}
\tabnote{Panel A uses equal observed-month averages over the 71 preperiod and 42 postperiod months; it reports the four underlying tail-by-age series before forming ratios.  The identity $-0.1581=-0.1098-0.0483$ closes to $2.5\times10^{-16}$.  Panel B's final three rows are an exact symmetric log-Shapley accounting of the same $-0.1581$ change.  All 88 family--tail--period denominator cells are classified before division: 52 are positive and 36 are structural absences, with no undefined active denominator.  No pseudocount is used.  These are model-free descriptions, not the grouped-binomial coefficient, confidence intervals, or causal shares.}
\end{table}


An outcome-free design audit explains why this accounting is not presented as
two saturated single-age regressions.  On the 52,884 single-age
occupation--month rows, unrestricted occupation-by-month indicators span the
entire row space and absorb all five requested exposure-by-post columns.  An
additive occupation-plus-calendar-month companion retains rank five for those
columns, but no outcome-bearing companion coefficients were estimated.  The
table therefore reports exact stock identities under the paper's normalization,
not a regression decomposition of the nonlinear grouped-binomial target.

\subsection{Leave-one-family and occupation influence}

Leave-one-family-out estimates are computed from the same family-conditioning program and reported with the omitted family's stock and target-information contribution.  These rows answer whether a broad comparison is concentrated; they are not used to choose which families belong in the main estimate.  Table~\ref{tab:app_lofo} displays the eight largest absolute movements.  Omitting office/administrative support or food preparation moves the coefficient in opposite directions by about 0.04 log point, and every paired deletion interval includes zero.

\begin{table}[!htbp]
\centering
\caption{Largest leave-one-family-out movements for the conditioned Q5 target}\label{tab:app_lofo}
\resizebox{\textwidth}{!}{%
\begin{tabular}{rlrrr}
\toprule
Omitted SOC2 & Family & Preperiod stock (\%) & Re-estimated coefficient & Movement [95\% paired CI] \\
\midrule
43 & Office and Administrative Support & 11.21 & -0.0614 & -0.0397 [-0.1131,0.0337] \\
35 & Food Preparation and Serving Related & 3.27 & 0.0174 & 0.0391 [-0.0292,0.1074] \\
31 & Healthcare Support & 3.66 & -0.0534 & -0.0317 [-0.0772,0.0139] \\
13 & Business and Financial Operations & 6.22 & 0.0099 & 0.0316 [-0.0006,0.0638] \\
39 & Personal Care and Service & 2.51 & -0.0457 & -0.0240 [-0.0496,0.0015] \\
25 & Educational Instruction and Library & 4.53 & -0.0351 & -0.0134 [-0.0383,0.0115] \\
17 & Architecture and Engineering & 2.28 & -0.0339 & -0.0123 [-0.0330,0.0085] \\
33 & Protective Service & 2.28 & -0.0096 & 0.0121 [-0.0050,0.0291] \\
\bottomrule
\end{tabular}}
\tabnote{The parent SOC2-by-calendar-month coefficient is $-0.0217$.  Families are ranked by the absolute coefficient movement, so this is an outcome-informed influence diagnostic rather than a preferred deletion rule.  Quintiles and the continuous centering scale remain fixed; all 22 deletions are available in the machine-readable file.}
\end{table}


Figure~\ref{fig:influential_between_family_paths} displays the economically
relevant endpoints behind those two largest leave-family-out movements: Q5
office/administrative-support stocks and Q1 food-preparation stocks.  These
families do not themselves span both tails, which is precisely why their paths
inform the broad between-family comparison rather than the sparse direct-tail
benchmark.  The selection is outcome-informed and the figure is diagnostic;
it neither identifies a pandemic mechanism nor supplies a deletion rule.

\begin{figure}[!htbp]
\centering
\includegraphics[width=0.96\textwidth]{r3_figure4_influential_family_paths.pdf}
\caption{Employment paths for the two largest leave-family-out movers}\label{fig:influential_between_family_paths}
\tabnote{The panels show Q1 food-preparation and serving occupations (SOC35) and Q5 office and administrative-support occupations (SOC43), the family-tail endpoints relevant to the broad Q5--Q1 comparison.  Young and older CPS-weighted employment stocks are indexed to their own 2019 means; October 2025 remains missing.  The families were selected after outcomes were observed because they generate the two largest absolute leave-family-out coefficient movements, so these paths are influence diagnostics rather than confirmatory evidence or a causal pandemic decomposition.  The red dashed line marks January 2023.}
\end{figure}

Figure~\ref{fig:family_paths} shows the four direct-tail families, selected by descending nuisance-adjusted Q5 information rather than coefficient sign.  Within each exposure tail, the young and older employment stocks are shown separately and indexed to their own 2019 means.  The figure therefore exposes which age series drives a tail's relative path without treating either age group as untreated.

\begin{figure}[!htbp]
\centering
\includegraphics[width=0.96\textwidth]{r3_figure3_family_paths.pdf}
\caption{Young and older employment stocks in direct-tail families}\label{fig:family_paths}
\tabnote{Families are selected by direct-tail, SOC2-by-calendar-month Q5 nuisance-adjusted information, with ties broken by SOC2 code; selection does not use coefficient sign.  Gray and blue distinguish Q1 and Q5; solid and dashed series distinguish ages 22--25 and 26--65.  Each series equals 100 at its own 2019 mean.  The symmetric-log vertical scale is linear from zero to 10 and logarithmic above 10 so that sparse-cell spikes do not obscure the other series; exact zero-stock months remain visible.  The red dashed line marks January 2023.  These are descriptive weighted-stock paths without design-based bands.}
\end{figure}

Occupation leave-one-out calculations are interpreted in the same way.  Fast-food and counter workers have high signed influence because low-exposure service employment recovered differently after the pandemic.  Other occupations push in the opposite direction.  Historical-contract diagnostics show that deleting the five largest absolute prior influences moves the coefficient toward zero, whereas deleting ten or twenty can make it more negative.  Symmetric influence trimming stays near the pooled estimate.  Because every such deletion is outcome adaptive, none replaces the broad-family regression.

Economically defined service exclusions are cleaner diagnostics but still change the population.  Removing Q1 food-service occupations attenuates the historical-contract absolute coefficient by less than one tenth; removing the broader SOC35/37/39 set of food, building-and-grounds, and personal-care service occupations does not attenuate it.  Protective-service occupations (SOC33) are not part of that exclusion.  These results reject a one-occupation explanation, not hypotheses about reopening, immigration, wages, or labor supply that the CPS design does not separately identify.
